\section{Positive-real maximal-modulus branch and endpoint asymptotics}\label{sec:positive-maximal-modulus-branch}

This section proves the two analytic inputs needed for the zeros: the positive real spectral root is uniquely dominant for $y>0$, and the square-root endpoint at $y=0$ has the cosine scaling used in Theorem~\ref{thm:endpoint-zero-quantization}.

\subsection{Positive branch parametrization and monotonicity}\label{subsec:positive-branch-parametrization}

\begin{lemma}[Local simplicity and persistence on $(0,\infty)$]\label{lem:local-simplicity-persistence}
Fix $y_0>0$, $y_0\ne1$.  The roots of $\Pi(\cdot,y_0)$ are simple.  Hence
\begin{enumerate}
\item\label{item:local-analyticity} each root extends analytically in $y$ near $y_0$;
\item\label{item:local-persistence} a strict modulus gap at $y_0$ persists locally.
\end{enumerate}
At $y_0=1$, the root $\lambda=2$ is simple and also persists analytically.
\end{lemma}
\begin{proof}
Corollary~\ref{cor:disc-factor} gives
\[
\Disc_\lambda\Pi=-y(y-1)(256y^3+411y^2+165y+32),
\]
and Lemma~\ref{lem:cubic-branch-real-root} places the unique real zero of the cubic at $\ySub<0$.  Thus the only positive discriminant zero is $y=1$.  Analyticity and gap persistence are the implicit-function theorem and continuity.  Finally,
\[
\Pi(\lambda,1)=(\lambda-2)(\lambda-1)(\lambda+1)^2,
\]
so $\lambda=2$ is simple.\label{rem:y1-reducible}
\end{proof}

\begin{proposition}[Preparatory positive-real parametrization]\label{prop:positive-maximal-modulus-branch}
For $\lambda>1$, set
\begin{equation}\label{eq:y-phys}
y_{\rm dom}(\lambda)=\lambda^2-\frac12-\frac12\sqrt{4\lambda^3-4\lambda+1}.
\end{equation}
Then $y_{\rm dom}$ maps $(1,\infty)$ bijectively onto $(0,\infty)$, with derivative
\begin{equation}\label{eq:dy-dlambda}
y_{\rm dom}'(\lambda)=2\lambda-\frac{3\lambda^2-1}{\sqrt{4\lambda^3-4\lambda+1}}>0.
\end{equation}
The inverse branch is denoted $\lambda_{\rm dom}(y)$.
\end{proposition}
\begin{proof}
The formula is the minus-sign sheet of \eqref{eq:y-from-E}.  Direct differentiation gives \eqref{eq:dy-dlambda}.  For $\lambda>1$,
\(F(\lambda)=4\lambda^3-4\lambda+1>0\) and \(3\lambda^2-1>0\), so positivity of \eqref{eq:dy-dlambda} is equivalent to
\[
2\lambda\sqrt{F(\lambda)}>3\lambda^2-1.
\]
Since both sides are positive, squaring is equivalent, and
\[
4\lambda^2F(\lambda)-(3\lambda^2-1)^2
=(\lambda-1)(\lambda+1)(16\lambda^3-9\lambda^2+1).
\]
For $\lambda>1$ the first two factors are positive, while
\(h(\lambda)=16\lambda^3-9\lambda^2+1\) satisfies \(h(1)=8>0\) and
\(h'(\lambda)=3\lambda(16\lambda-6)>0\).  Hence \(y'_{\rm dom}(\lambda)>0\).  The limits are $0$ as $\lambda\downarrow1$ and $+\infty$ as $\lambda\to\infty$, so the map is bijective.
\end{proof}

\begin{lemma}[Comparison of the two positive branches for $y>1$]\label{lem:ygt1-sheet-comparison}
For $\lambda>1$, define
\[
y_-(\lambda)=\lambda^2-\frac12-\frac12\sqrt{4\lambda^3-4\lambda+1},
\qquad
y_+(\lambda)=\lambda^2-\frac12+\frac12\sqrt{4\lambda^3-4\lambda+1}.
\]
Then $y_+$ is strictly increasing and $y_+>y_-$.  If $y=y_-(q)=y_+(p)$ with $p,q>1$, then $q>p$.
\end{lemma}
\begin{proof}
Here $y_+'>0$ because $2\lambda+(3\lambda^2-1)/\sqrt{4\lambda^3-4\lambda+1}>0$.  Thus $y_+(q)>y_-(q)=y=y_+(p)$ implies $q>p$.
\end{proof}

\begin{lemma}[No small positive roots for $y>1$]\label{lem:ygt1-no-small-positive-roots}
For every $y>1$ and every $x\in[0,1]$, $\Pi(x,y)>0$.  Hence every positive real root of $\Pi(\cdot,y)$ lies in $(1,\infty)$.
\end{lemma}
\begin{proof}
As a quadratic in $y$,
\[
\Pi(x,y)=y^2+(1-2x^2)y+x^4-x^3-x^2+x.
\]
For $x\in[0,1]$ and $y>1$, $\partial_y\Pi=2y+1-2x^2>0$.  Moreover
\[
\Pi(x,1)=(1-x)(2-x)(1+x)^2\ge0.
\]
Therefore $\Pi(x,y)>\Pi(x,1)\ge0$.
\end{proof}

\begin{lemma}[Primitive Sturm count on the positive axis]\label{lem:positive-axis-sturm-count}
Let $y_0=(17+\sqrt{297})/2$, the positive root of $y^2-17y-2$.  The primitive rational Sturm chain for \(\Pi(\lambda,y)\) is
\begin{align*}
S_0&=\Pi(\lambda,y),\\
S_1&=\partial_\lambda\Pi(\lambda,y),\\
16S_2&=(16y+11)\lambda^2+(4y-10)\lambda-16y^2-16y-1,\\
-\frac{(16y+11)^2}{16}S_3
&=4\lambda y^2-68\lambda y-8\lambda-64y^3-41y^2+25y+8,\\
S_4&=-\frac{y(y-1)(16y+11)^2(256y^3+411y^2+165y+32)}{256(y^2-17y-2)^2}.
\end{align*}
On each cell avoiding the denominator wall \(y^2-17y-2=0\), the divisions in \(\QQ(y)[\lambda]\) satisfy
\[
\operatorname{rem}(S_0,S_1)=-S_2,
\qquad
\operatorname{rem}(S_1,S_2)=-S_3,
\qquad
\operatorname{rem}(S_2,S_3)=-S_4.
\]
The sign-controlling factors are
\begin{equation}\label{eq:positive-axis-sturm-sign-factors}
y,\quad y-1,\quad 256y^3+411y^2+165y+32,
\quad y^2-17y-2.
\end{equation}
The Sturm divisions are read on cells avoiding the denominator wall \(y_0\).  On the positive axis the corrected variation table is
\begin{equation}\label{eq:positive-axis-sturm-sign-table-main}
\begin{array}{c|c|c|c}
\text{parameter domain} & \text{signs at }-\infty & \text{signs at }+\infty & \text{variation drop}\\
\hline
0<y<1 & (+,-,+,-,+) & (+,+,+,+,+) & 4\\
1<y<y_0 & (+,-,+,-,-) & (+,+,+,+,-) & 2\\
y>y_0 & (+,-,+,+,-) & (+,+,+,-,-) & 2.
\end{array}
\end{equation}
Consequently \(\Pi(\cdot,y)\) has four real roots for \(0<y<1\) and two real roots for \(y>1\).
\end{lemma}

\begin{proof}
The displayed primitive rational chain has only one positive denominator wall, the positive root \(y_0\) of \(y^2-17y-2\).  Its leading signs at \(\pm\infty\), read from the displayed formulae for \(S_0,\ldots,S_4\), are exactly those in \eqref{eq:positive-axis-sturm-sign-table-main}.  Since \(y_0\ne1\) and the cubic factor in the discriminant
\[
\Disc_\lambda\Pi=-y(y-1)(256y^3+411y^2+165y+32)
\]
is positive for \(y_0>0\), the discriminant is nonzero at \(y_0\).  Thus no root can appear, disappear, or collide at the Sturm-chain denominator wall, and the two adjacent \(y>1\) cells have the same real-root count.  This proves the stated counts.
\end{proof}

\begin{lemma}[No negative roots for $y>1$]\label{lem:ygt1-no-negative-roots}
For every $y>1$, the polynomial $\Pi(\cdot,y)$ has no negative real root.
\end{lemma}
\begin{proof}
For $y>1$, Proposition~\ref{prop:positive-maximal-modulus-branch} gives one positive root
$q>1$ on the sheet $y=y_-(q)=y_{\rm dom}(q)$.  Since $y_+$ is strictly increasing,
$y_+(1)=1$, and $y_+(\lambda)\to\infty$ as $\lambda\to\infty$, Lemma~\ref{lem:ygt1-sheet-comparison}
gives a second positive root $p>1$ on the sheet $y=y_+(p)$.  The two sheet equations are precisely
the two solutions in $y$ obtained from $\Pi(\lambda,y)=0$, so both $p$ and $q$ are roots of
$\Pi(\cdot,y)$.  The preceding variation table in Lemma~\ref{lem:positive-axis-sturm-count} gives
exactly two real roots for every $y>1$.  Since two positive real roots have already been produced,
there is no remaining real root, in particular no negative real root.
\end{proof}

\subsection{Algebraic dominance criterion on the positive axis}\label{subsec:positive-axis-dominance}

\begin{lemma}[Positive-axis dominance criterion]\label{lem:positive-axis-dominance-criterion}
For $y>0$, the following hold.
\begin{enumerate}
\item The discriminant and endpoint evaluations are
\begin{equation}\label{eq:positive-axis-discriminant}
\Disc_\lambda\Pi=-y(y-1)(256y^3+411y^2+165y+32),
\end{equation}
\begin{equation}\label{eq:positive-axis-resultants}
\Pi(0,y)=y(y+1),\qquad \Pi(1,y)=\Pi(-1,y)=y(y-1).
\end{equation}
\item The Sturm count is the earlier Lemma~\ref{lem:positive-axis-sturm-count}: the denominator wall is the positive root \(y_0=(17+\sqrt{297})/2\) of the corrected factor \(y^2-17y-2\), the variation table is \eqref{eq:positive-axis-sturm-sign-table-main}, and the resulting real-root counts are four for \(0<y<1\) and two for \(y>1\).
\item Opposite real roots can have equal modulus only at $y=0$ or $y=1$:
\begin{equation}\label{eq:opposite-modulus-resultant}
\Pi(-x,y)-\Pi(x,y)=2x(x^2-1),\qquad
\Res_x(\Pi(x,y),2x(x^2-1))=16y^3(y-1)^2(y+1).
\end{equation}
\item If $0<y<1$, the four real roots satisfy
\begin{equation}\label{eq:positive-axis-root-boxes}
r_1<-1<r_2<0<r_3<1<r_4,
\end{equation}
and $r_4=\lambda_{\rm dom}(y)$ has strictly largest modulus.  If $y>1$, the two real roots are positive and lie in $(1,\infty)$; writing them $p<q$, the larger one is $q=\lambda_{\rm dom}(y)$, and the nonreal pair has modulus $<q$.
\end{enumerate}
\end{lemma}
\begin{proof}
 The discriminant is Corollary~\ref{cor:disc-factor}, the endpoint evaluations are substitutions, and the Sturm count is Lemma~\ref{lem:positive-axis-sturm-count}.

For $0<y<1$, \eqref{eq:positive-axis-resultants} and the leading coefficient place one real root in each interval of \eqref{eq:positive-axis-root-boxes}; the Sturm count shows these are all roots.  The roots in $(-1,0)$ and $(0,1)$ have smaller modulus than $r_4$ by inspection of the intervals.  It remains to orient the comparison between $r_4$ and the root $r_1<-1$.  On $(0,1)$ the discriminant is nonzero, so $r_1(y)$ and $r_4(y)$ vary continuously.  Opposite real roots can have equal modulus only when the resultant in \eqref{eq:opposite-modulus-resultant} vanishes, which is impossible for $0<y<1$; hence the sign of $r_4(y)+r_1(y)$ is constant on this interval.  To determine it, use the factorization at $y=1$:
\[
\Pi\!\left(-\frac32,1\right)=\frac{35}{16}>0,
\qquad
\Pi\!\left(\frac32,1\right)=-\frac{25}{16}<0.
\]
For all $y<1$ sufficiently close to $1$, continuity gives $\Pi(-3/2,y)>0$ and $\Pi(3/2,y)<0$, while $\Pi(-1,y)=\Pi(1,y)=y(y-1)<0$.  Since there is exactly one root in $(-\infty,-1)$ and exactly one root in $(1,\infty)$, this places $r_1\in(-3/2,-1)$ and $r_4>3/2$ for such $y$.  Thus $r_4>|r_1|$ near $1$, and the no-equality statement propagates this strict inequality throughout $0<y<1$.

For $y>1$, Proposition~\ref{prop:positive-maximal-modulus-branch} gives one positive root
$q>1$ on the sheet $y=y_-(q)=y_{\rm dom}(q)$.  Since $y_+$ is strictly increasing,
$y_+(1)=1$, and $y_+(\lambda)\to\infty$ as $\lambda\to\infty$, there is also a unique
$p>1$ with $y=y_+(p)$.  The two sheet equations are precisely the two solutions in $y$
obtained from $\Pi(\lambda,y)=0$, so $p$ and $q$ are roots of $\Pi(\cdot,y)$.  The same Sturm
variation table in \eqref{eq:positive-axis-sturm-sign-table-main} gives exactly two real
roots for every $y>1$; since $p$ and $q$ are already positive real roots, no negative real root remains.
By Lemma~\ref{lem:ygt1-no-small-positive-roots} both positive
roots lie in $(1,\infty)$.  Lemma~\ref{lem:ygt1-sheet-comparison} gives $q>p$, so writing
the two real roots as $1<p<q$ identifies the larger root with the minus sheet
$\lambda_{\rm dom}$.

It remains only to compare $q$ with the nonreal pair.  Let that pair be $z,\overline z$ and put $r=|z|^2$.  Since the coefficient of $\lambda^3$ is $-1$, we have
\[
z+\overline z=1-p-q.
\]
Vieta's formula for the coefficient of $\lambda$ gives
\[
-\bigl(pq(z+\overline z)+r(p+q)\bigr)=1,
\]
or equivalently
\[
r=\frac{pq(p+q-1)-1}{p+q}.
\]
Therefore
\[
q^2-r=
\frac{q^2(p+q)-pq(p+q-1)+1}{p+q}
=\frac{q(q-p)(p+q)+pq+1}{p+q}>0,
\]
because $q>p>1$.  Thus $|z|^2=r<q^2$, so $|z|<q$.
\end{proof}

\begin{theorem}[Canonical positive-axis dominance]\label{thm:canonical-branch-dominance}\label{eq:consolidated-disc-factor}
For every $y>0$, the branch $\lambda_{\rm dom}(y)>1$ of Proposition~\ref{prop:positive-maximal-modulus-branch} is a root of $\Pi(\cdot,y)$ and is the unique root of maximal modulus.  The branch-value fibre counts used at $y=1$ are those of Proposition~\ref{prop:branch-value-real-root-counts}; at $y=1$, this branch gives the surviving simple spectral root $\lambda_{\rm dom}(1)=2$.
\end{theorem}
\begin{proof}
For $y\ne1$ this is Lemma~\ref{lem:positive-axis-dominance-criterion}, whose Sturm table uses the denominator wall given by the positive root of the corrected factor $y^2-17y-2$.  At $y=1$ the factorization $(\lambda-2)(\lambda-1)(\lambda+1)^2$ gives the unique largest modulus root $2$, and the reduced recurrence keeps precisely this surviving factor by Proposition~\ref{prop:minimal-order}.
\end{proof}

\begin{remark}[Spectral-radius reading on the positive axis]\label{cor:positive-axis-spectral-radius-reduction}\label{cor:normalized-branch-classification}\label{eq:carried-normalized-branch-change}\label{prop:positive-side-real-scaling-stability}
For $y>0$, Theorem~\ref{thm:canonical-branch-dominance} makes $\lambda_{\rm dom}(y)$ the spectral radius of the active transfer spectrum.  Under a real weighted pullback of the canonical member, the corresponding statement is used only through Lemma~\ref{lem:real-weighted-pullback-dominance}: the spectral coordinate is multiplied by a common nonzero real factor, while the shifted parameter is first returned to the canonical positive axis.  No dominance comparison is transported through a horizontal translation in $\lambda$, and no separate parameter-space dominance theorem is used.
\end{remark}

\subsection{Endpoint at \texorpdfstring{$y=0$}{y=0} and Puiseux expansion}\label{subsec:endpoint-puiseux}

\begin{proposition}[Endpoint Puiseux expansion at $y=0$]\label{prop:endpoint-puiseux}
This is the local $\Pi$-specific expansion used by the endpoint zero theorem; no general endpoint machinery is invoked in the main proof.  With $y=s^2$ near $s=0$, the two branches colliding at $\lambda=1$ satisfy
\begin{equation}\label{eq:y-puiseux}
y=s^2,
\end{equation}
\begin{equation}\label{eq:lambda-puiseux}
\lambda_\pm(s)=1\pm\frac{s}{\sqrt2}+\frac58s^2\mp\frac{43}{64\sqrt2}s^3+O(s^4),
\end{equation}
and
\begin{equation}\label{eq:f-puiseux}
\log\lambda_\pm(s)=\pm\frac{s}{\sqrt2}+\frac38s^2\mp\frac{217}{192\sqrt2}s^3+O(s^4).
\end{equation}
The local density coefficient at the endpoint is
\begin{equation}\label{eq:rho-endpoint}
\rho(0)=\frac1{\pi\sqrt2}.
\end{equation}
\end{proposition}
\begin{proof}
Substitute $y=s^2$ and solve $\Pi(\lambda,s^2)=0$ by formal power series at $\lambda=1$.  The logarithmic expansion follows by expanding $\log(1+(\lambda_\pm-1))$.  Appendix~\ref{app:endpoint-branching} records the coefficient computation and residue normalization.  The displayed cubic logarithmic term contributes only $O_K(m^{-2})$ to $m\log\lambda_\pm(iu/m)$ on fixed compact endpoint scales, so the fixed-window zero theorem below uses only the coefficients $a=1/\sqrt2$ and $b=3/8$ to first correction order.
\end{proof}

\subsection{Endpoint cosine law and Rouch\'e count}\label{subsec:endpoint-quantization}

\begin{lemma}[Removable endpoint amplitudes]\label{lem:endpoint-amplitude-removability}
Let $y=s^2$ and let $\lambda_\pm(s)$ be the two branches through $\lambda=1$ in \eqref{eq:lambda-puiseux}.  On the branch-coordinate disc $|s|<1/2$, the corresponding partial-fraction amplitudes
\[
C_\pm(s)=\frac{\lambda_\pm(s)^3N(1/\lambda_\pm(s),s^2)}{\partial_\lambda\Pi(\lambda_\pm(s),s^2)}
\]
have removable singularities at $s=0$ and extend holomorphically there with
\begin{equation}\label{eq:endpoint-amplitude-removable-main}
C_\pm(s)=\frac12\mp\frac{3\sqrt2}{16}s+O(s^2).
\end{equation}
Consequently there is a fixed radius $\rho_{\rm amp}$, $0<\rho_{\rm amp}\le1/2$, such that both extended amplitudes are nonzero on $|s|<\rho_{\rm amp}$; we choose $\rho_{\rm amp}$ so that $|C_\pm(s)-1/2|<1/4$ on this disc.
\end{lemma}
\begin{proof}
Away from $s=0$ the displayed residue formula is the usual simple-root partial-fraction coefficient.  Appendix~\ref{app:endpoint-branching}, especially Lemmas~\ref{lem:endpoint-coefficient-computation} and~\ref{lem:plus-branch-residue-cancellation}, computes the Taylor numerator and denominator along the plus branch as
\[
\lambda_+(s)^3N(1/\lambda_+(s),s^2)=\sqrt2\,s+\frac34s^2+O(s^3),
\qquad
\partial_\lambda\Pi(\lambda_+(s),s^2)=2\sqrt2\,s+3s^2+O(s^3).
\]
The common linear zero cancels, and division gives $C_+(s)=1/2-(3\sqrt2/16)s+O(s^2)$.  The branch symmetry $\lambda_-(s)=\lambda_+(-s)$ gives the minus-branch expansion.  Thus the singularity at $s=0$ is removable and \eqref{eq:endpoint-amplitude-removable-main} holds.  The nonvanishing radius follows from continuity of the holomorphic extensions and the value $C_\pm(0)=1/2$.
\end{proof}

Throughout the endpoint-window arguments, the compact set \(K\subset\CC\), the positive window size \(U\), and any displayed index \(j\), \(k\), or finite set of indices are fixed before \(m\to\infty\).  All constants implicit in \(O_K(\cdot)\), \(O_U(\cdot)\), and \(O_k(\cdot)\) may depend on these fixed choices, but never on \(m\).  No assertion in Proposition~\ref{prop:fixed-window-endpoint-isolation} or Theorem~\ref{thm:endpoint-zero-quantization} is uniform for growing windows or growing indices.

\begin{proposition}[Local endpoint transfer and fixed-window isolation for $\Pi$]\label{prop:fixed-window-endpoint-isolation}
The comparison function is the shifted cosine $G_m$ used below; the fixed-window convention in the preceding paragraph applies to this proposition.
This proposition is the compact local endpoint-transfer statement for the canonical quartic.  Use the endpoint constants computed in Appendix~\ref{app:endpoint-branching},
\begin{equation}\label{eq:main-endpoint-constant-handoff}
a=\frac1{\sqrt2},\qquad b=\frac38,
\qquad A_*=\frac12,
\qquad B_*=\frac{3\sqrt2}{16},
\end{equation}
where $a$ is the linear logarithmic branch coefficient, $b$ is the quadratic logarithmic coefficient, and $A_*,B_*$ are the constant and odd linear terms of the two colliding partial-fraction amplitudes from Lemma~\ref{lem:endpoint-amplitude-removability}.
Define
\[
F_m(u):=e^{3u^2/(8m)}Z_m\!\left(-\frac{u^2}{m^2}\right).
\]
For every compact $K\subset\CC$,
\begin{equation}\label{eq:quartic-corollary-cosine}
F_m(u)=\cos\!\left(\left(a-\frac{B_*}{A_*m}\right)u\right)+O_K(m^{-2}),
\end{equation}
which is equivalently
\begin{equation}\label{eq:endpoint-scaling-normalized}
Z_m\!\left(-\frac{u^2}{m^2}\right)=
\exp\!\left(-\frac{3u^2}{8m}\right)
\left[\cos\!\left(\frac{u}{\sqrt2}-\frac{3u}{4\sqrt2\,m}\right)+O_K(m^{-2})\right].
\end{equation}
If $U>0$ is not one of
$\nu_j^{(0)}=(2j-1)\pi/\sqrt2$, and $N(U)=\#\{j\ge1:\nu_j^{(0)}<U\}$, then for all sufficiently large $m$ the rectangle
\[
Q_U=\{u\in\CC:0\le \Re u\le U,
\ |\Im u|\le1\}
\]
contains exactly $N(U)$ zeros of $u\mapsto Z_m(-u^2/m^2)$, all simple and real positive, one near each $\nu_j^{(0)}<U$, and no others.  For fixed $j$,
\begin{equation}\label{eq:quartic-u-location}
u_{m,j}=\frac{\pi(2j-1)}{\sqrt2}\left(1+\frac{3}{4m}\right)+O_j(m^{-2}),
\end{equation}
and
\begin{equation}\label{eq:quartic-y-location-local}
y_{m,j}^{(U)}=-\frac{u_{m,j}^2}{m^2}
=-\frac{\pi^2(2j-1)^2}{2m^2}-\frac{3\pi^2(2j-1)^2}{4m^3}+O_j(m^{-4}).
\end{equation}
\end{proposition}
\begin{proof}
Appendix~\ref{app:endpoint-branching} is used only through the fixed local data in Proposition~\ref{prop:endpoint-expansion}.  With $y=s^2$ and $s=iu/m$, that proposition gives the two colliding branches and amplitudes in the common branch coordinate
\[
\log\lambda_\pm(s)=\pm as+bs^2+O(s^3),
\qquad
C_\pm(s)=A_*\mp B_*s+O(s^2),
\]
with the constants in \eqref{eq:main-endpoint-constant-handoff}, and it gives the combined contribution of the two noncolliding branches as $O_K(m^{-2})$ after the same substitution.  Hence, uniformly for $u\in K$,
\[
\lambda_\pm(iu/m)^m
=
\exp\!\left(\pm iau-\frac{bu^2}{m}+O_K(m^{-2})\right),
\qquad
C_\pm(iu/m)=A_*\mp iB_*\frac{u}{m}+O_K(m^{-2}).
\]
Adding the two colliding contributions and absorbing the noncolliding remainder gives
\begin{align*}
Z_m\!\left(-\frac{u^2}{m^2}\right)
&=
e^{-bu^2/m}
\left[
2A_*\cos(au)+\frac{2B_*u}{m}\sin(au)+O_K(m^{-2})
\right]  \\
&=
e^{-3u^2/(8m)}
\left[
\cos\!\left(\frac{u}{\sqrt2}\right)+\frac{3u}{4\sqrt2\,m}\sin\!\left(\frac{u}{\sqrt2}\right)+O_K(m^{-2})
\right],
\end{align*}
because $2A_*=1$ and $2B_*=3\sqrt2/8$.  On the fixed compact $K$, the elementary identity
\[
\cos(\theta-\delta)=\cos\theta+\delta\sin\theta+O_K(m^{-2}),
\qquad
\theta=\frac{u}{\sqrt2},
\quad
\delta=\frac{3u}{4\sqrt2\,m},
\]
is exactly \eqref{eq:quartic-corollary-cosine}, and multiplying by $e^{-3u^2/(8m)}$ gives \eqref{eq:endpoint-scaling-normalized}.

Now fix $U$ before taking $m\to\infty$, and assume $U\ne\nu_j^{(0)}$ for all $j$.  Put
\[
G_m(u)=\cos\!\left(\left(a-\frac{B_*}{A_*m}\right)u\right),
\]
so that \eqref{eq:quartic-corollary-cosine} reads \(F_m=G_m+O_K(m^{-2})\) on each fixed compact set.  Let
\[
\nu_j^{(m)}=\frac{(2j-1)\pi/2}{a-B_*/(A_*m)}
\]
be the real simple zeros of \(G_m\).  Since $U$ is not a limiting threshold, choose $\eta>0$ so that, for all sufficiently large $m$, the closed real-centred discs
\[
D_j^{\rm fix}=\{u\in\CC:|u-\nu_j^{(m)}|\le\eta\},
\qquad \nu_j^{(0)}<U,
\]
are pairwise disjoint, conjugation-invariant, contained in $Q_U$, and avoid $\partial Q_U$.  Let
\[
\Omega_U^{\rm fix}=Q_U\setminus\bigcup_{\nu_j^{(0)}<U}\operatorname{int}D_j^{\rm fix}.
\]
The boundary of \(\Omega_U^{\rm fix}\) is piecewise smooth and contains no zero of \(G_m\).  On the fixed compact set \(Q_U\), write
\[
F_m(u)=G_m(u)+H_m(u),
\qquad |H_m(u)|\le C_Um^{-2}
\]
for all sufficiently large $m$.

We first isolate on the fixed discs only to obtain the count.  On the compact set consisting of \(\Omega_U^{\rm fix}\) and the finitely many circle boundaries one has \(|G_m|\ge c_U>0\) for all sufficiently large $m$.  Thus \(|H_m|<|G_m|\) there.  Rouch\'e's theorem on each \(D_j^{\rm fix}\) gives exactly one zero of \(F_m\) in \(D_j^{\rm fix}\), with multiplicity one.  Applying the argument-principle form of Rouch\'e's theorem to the multiply connected domain \(\Omega_U^{\rm fix}\) gives the same number of zeros for \(F_m\) and \(G_m\) in \(\Omega_U^{\rm fix}\), namely none.

The fixed-disc count alone is not used for the location estimate.  We now shrink to the natural perturbative scale.  For a fixed counted index $j$, put
\[
D_{j,m}(R)=\{u\in\CC:|u-\nu_j^{(m)}|\le Rm^{-2}\}.
\]
Since the zeros of $G_m$ are simple and $a-B_*/(A_*m)$ stays in a compact subinterval of $(0,\infty)$, Taylor's formula at $\nu_j^{(m)}$ gives constants $c_j>0$ and $R_j^0>0$ such that, whenever $R\le R_j^0m^2$ and $m$ is large,
\begin{equation}\label{eq:quantitative-cosine-lower-bound}
|G_m(u)|\ge c_j Rm^{-2}
\qquad\text{on } |u-\nu_j^{(m)}|=Rm^{-2}.
\end{equation}
Choose $R_j$ so large that $c_jR_j>C_U$.  Then \(|H_m|<|G_m|\) on \(\partial D_{j,m}(R_j)\), and Rouch\'e's theorem gives one zero of $F_m$ inside \(D_{j,m}(R_j)\).  The fixed-disc count has already shown that there is only one zero in \(D_j^{\rm fix}\), so the zero counted above is the same one.  Hence
\begin{equation}\label{eq:quantitative-u-location}
u_{m,j}=\nu_j^{(m)}+O_j(m^{-2}).
\end{equation}

For completeness, the same estimate also excludes hidden zeros in the complement of the $m^{-2}$ discs.  On the complement of the finitely many interiors \(D_{j,m}(R_j)\) inside $Q_U$, the only possible small values of $G_m$ occur near its already listed simple zeros; the lower bound \eqref{eq:quantitative-cosine-lower-bound}, after increasing the constants for the finite set of indices, gives \(|G_m|>|H_m|\) on each new circle boundary, while the outer boundary has a fixed positive lower bound because $U$ is not a threshold.  Rouch\'e on the resulting multiply connected domain therefore gives no further zeros.

The polynomial $Z_m(y)$ has real coefficients, so the function \(u\mapsto Z_m(-u^2/m^2)\), and hence also \(F_m\), satisfies \(F_m(\overline u)=\overline{F_m(u)}\).  Since each $D_j^{\rm fix}$ is invariant under conjugation and contains exactly one zero, that zero is fixed by conjugation and hence real.  The zeros are simple because the local Rouch\'e count preserves the multiplicity-one count of the shifted cosine zeros.  At every counted zero one has $u>0$, so the map $u\mapsto y=-u^2/m^2$ is locally biholomorphic there, with $dy/du=-2u/m^2\ne0$; the conjugate zero at $-u$ of the even function $Z_m(-u^2/m^2)$ represents the same $y$-zero and lies outside the chosen positive $u$-rectangle.

The center of the shifted cosine satisfies
\[
\left(\frac1{\sqrt2}-\frac{3}{4\sqrt2\,m}\right)u_{m,j}
=
\frac{(2j-1)\pi}{2}+O_j(m^{-2}).
\]
Inverting the fixed scalar factor gives \eqref{eq:quartic-u-location}, and substituting $y=-u^2/m^2$ gives \eqref{eq:quartic-y-location-local}.\label{par:endpoint-theorem-inputs}
\end{proof}

\begin{lemma}[Endpoint rectangle boundary and $u$--$y$ transfer]\label{lem:endpoint-boundary-transfer}
Fix $U>0$ with $U\ne (2j-1)\pi/\sqrt2$ for all $j\ge1$, and let $Q_U$ be the rectangle in Proposition~\ref{prop:fixed-window-endpoint-isolation}.  For all sufficiently large $m$, the function $u\mapsto Z_m(-u^2/m^2)$ has no zero on $\partial Q_U$.  Its zeros in $Q_U$ are exactly the positive real zeros supplied by Proposition~\ref{prop:fixed-window-endpoint-isolation}, and the map
\[
u\longmapsto y=-\frac{u^2}{m^2}
\]
transfers them bijectively, with multiplicity preserved, to the zeros of $Z_m(y)$ in $[-U^2/m^2,0)$.
\end{lemma}
\begin{proof}
Let $a_m=a-B_*/(A_*m)$.  The comparison function $G_m(u)=\cos(a_mu)$ has only real zeros, at $\nu_j^{(m)}=(2j-1)\pi/(2a_m)$.  Since $U$ is not a limiting threshold, no zero of $G_m$ lies on $\partial Q_U$ for all sufficiently large $m$: $G_m(0)=1$, the horizontal sides have nonzero imaginary part, and the right side $\Re u=U$ stays a fixed positive distance from the finite set of real zeros in the window.  Hence $|G_m|$ has a positive lower bound on $\partial Q_U$ for large $m$.  The uniform estimate $F_m=G_m+O_U(m^{-2})$ and the nonvanishing factor $e^{3u^2/(8m)}$ then exclude zeros of $u\mapsto Z_m(-u^2/m^2)$ on $\partial Q_U$.

Proposition~\ref{prop:fixed-window-endpoint-isolation} gives exactly one simple positive real zero $u_{m,j}$ for each threshold below $U$ and no other zero in $Q_U$.  The map $u\mapsto -u^2/m^2$ is injective on $0<u<U$ and maps this interval onto $(-U^2/m^2,0)$.  At each counted zero, $dy/du=-2u/m^2\ne0$, so the holomorphic change of variable preserves multiplicity.  Conversely, any negative zero $y\in(-U^2/m^2,0)$ has the unique positive preimage $u=m\sqrt{-y}\in(0,U)$ in $Q_U$, so no zero in the $y$-window is missed.
\end{proof}

\begin{definition}[Endpoint zero ordering]\label{def:endpoint-zero-ordering}
In a real negative endpoint window, zeros are ordered from $0$ leftwards; equivalently $y_{m,k}=-u_{m,k}^2/m^2$ with $0<u_{m,1}<u_{m,2}<\cdots$.
\end{definition}

\begin{theorem}[Fixed-index, fixed-window endpoint zero law]\label{thm:endpoint-zero-quantization}
For the canonical quartic,
\begin{equation}\label{eq:endpoint-scaling}
Z_m\!\left(-\frac{u^2}{m^2}\right)=
\exp\!\left(-\frac{3u^2}{8m}\right)
\left[
\cos\!\left(\frac{u}{\sqrt2}-\frac{3u}{4\sqrt2\,m}\right)+O_K(m^{-2})
\right]
\end{equation}
for every fixed compact $K\subset\CC$.  If $U>0$ is fixed and $U\ne(2j-1)\pi/\sqrt2$ for all $j\ge1$, then for all sufficiently large $m$ the window $[-U^2/m^2,0)$ contains exactly
\[
N(U)=\#\left\{j\ge1:\frac{(2j-1)\pi}{\sqrt2}<U\right\}
=\left\lfloor\frac{U}{\pi\sqrt2}+\frac12\right\rfloor
\]
negative real zeros of $Z_m$, all simple and ordered by Definition~\ref{def:endpoint-zero-ordering}.  For each fixed $k$ in such a window,
\begin{equation}\label{eq:endpoint-zero-quantization}
y_{m,k}= -\frac{\pi^2(2k-1)^2}{2m^2}-\frac{3\pi^2(2k-1)^2}{4m^3}+O_k(m^{-4}).
\end{equation}
Here $K$, $U$, and any finite set of indices are fixed before $m\to\infty$; the theorem makes no assertion for growing endpoint windows or growing $k$.  The proof uses only the local endpoint-transfer Proposition~\ref{prop:fixed-window-endpoint-isolation} after the constants have been computed.
\end{theorem}
\begin{proof}
The estimate \eqref{eq:endpoint-scaling} is Proposition~\ref{prop:fixed-window-endpoint-isolation}.  The same proposition gives the Rouch\'e count in $Q_U$ and the locations \eqref{eq:quartic-u-location}.  Lemma~\ref{lem:endpoint-boundary-transfer} excludes boundary zeros and transfers the positive $u$-zeros bijectively, with multiplicity preserved, to the negative $y$-zeros in the window.  Squaring \eqref{eq:quartic-u-location} gives \eqref{eq:endpoint-zero-quantization}.
\end{proof}
